% \chapter{Travaux à réaliser pour les six prochaines semaines}

% Les tâches prévues et restantes pour l'aboutissement du projet sont réparties selon les axes suivants :

% Pour la régulation PI :
% \begin{itemize}
%     \item Déterminer expérimentalement la valeur optimale du gain $K_{pi}$.
% \end{itemize}

% Pour la régulation par retour d'état :
% \begin{itemize}
%     \item Identifier un placement de pôles assurant la robustesse et la stabilité de la régulation.
%     \item Analyser les simulations afin de définir un pôle dominant optimal.
% \end{itemize}

% Pour la régulation PID :
% \begin{itemize}
%     \item Réaliser une simulation complète du système.
%     \item Vérifier la validité de la synthèse théorique.
%     \item Appliquer les paramètres sur un seul inducteur pour identifier le comportement réel.
%     \item Étendre l'application de la commande aux quatre inducteurs.
%     \item Déterminer les pôles du système en boucle fermée.
%     \item Transposer les pôles identifiés sur la régulation d'état.
% \end{itemize}

% Pour le programme principal du DSP :
% \begin{itemize}
%     \item Intégrer une procédure d'atterrissage en douceur pour la maquette.
%     \item Implémenter une réinitialisation systématique des variables d'état.
%     \item Optimiser les conditions et la logique des états 0 et 8 de la machine d'états.
%     \item Mettre en place un système de diagnostic et de remonte des messages d'erreur.
% \end{itemize}

% Pour le microcontrôleur ESP32 et l'interface web :
% \begin{itemize}
%     \item Intégrer l'affichage en temps réel de la mesure de courant.
%     \item Corriger le dysfonctionnement lié au chargement infini de la page web.
%     \item Générer un code QR pour faciliter la connexion directe au portail web.
%     \item Ajouter des éléments d'interface interactifs (ex. : curseurs/sliders) pour modifier dynamiquement les coefficients de régulation.
% \end{itemize}

% De manière générale :
% \begin{itemize}
%     \item Continuer de chercher des problèmes
%     \item Se poser mille et une question sur pourquoi ça ne marche pas
% \end{itemize}
